% !TeX root = ../icml2026.tex
\section{Introduction}

Embedding-based retrieval systems answer queries by vector comparison. One common setting of such systems is that the system maintains each element $x_i\in X$ as a vector $\vx_i\in \real^d$, where $X$ is the universe set of elements and $\real^d$ is the vector space\footnote{Other metric spaces may apply. $\real^d$ is chosen because it is the most commonly used.}.
Each query $q$ to answer is also embedded as a vector $\vw_q\in \real^d$. The answers to a query $q$ are retrieved by (1) comparing $\vw_q$ against each of the $\vx_i$ with a scoring function $s(\cdot, \cdot)$ to obtain a score $s_{i|q} = s(\vx_i, \vw_q)$ and (2) retrieving the answers with the $k$ largest scores $s_{i|q}$\footnote{Other notions of answer sets may apply, but top-$k$ is chosen also because its common usage.}.
This paper investigates a fundamental question of the aforementioned retrieval systems. One informal statement is:
\begin{tcolorbox}[colback=yellow!20, colframe=black, sharp corners, boxrule=1pt]
Given the universe $X$ with \underline{$m$ elements} and a \underline{scoring function $s$}, what is the \textbf{minimal} \underline{dimension $d$} of $\real^d$ such that \textbf{every} query with \underline{at most $k$ answers} is perfectly retrievable?
\end{tcolorbox}
This question concerns the ``Minimal Embeddable Dimension (MED)'' $d$, as will be formally defined in Section~\ref{sec:problem-definition}, which depends on the function $s$, the cardinality $m$ of the universe $X$, and also the cardinality $k$ of the answer set to queries we are interested in. For convenience, we say a set of embeddings for elements and subsets is an \textbf{embeddable configuration} if and only if all answers can be retrieved by score comparison.

The study of the minimal embeddable dimension is fundamentally related to several classic research directions, including the VC-dimension in statistical learning theory~\citep{mohri2018foundations} and the $k$-set problem in combinatorial geometry~\citep{matousek2013lectures}. It is widely acknowledged that this question is related to the fundamental boundary underlying modern information retrieval~\citep{lee2019latent,weller2025mfollowir}, data compression~\citep{andoni2008near}, and representation learning~\citep{izacard2021unsupervised,wang2022text}.


The effectiveness of embedding-based retrieval has been demonstrated to depend on the dimension $d$ of the space when the size $m$ of the universe is large~\citep{yin2018dimensionality,reimers2021curse}. Very recently, \citet{weller2025theoretical} suggested that the fundamental property of the underlying geometric space limits the effectiveness of embedding-based retrieval, as supported by both theoretical findings and numerical simulations. On the theory side, they considered a retrieval problem defined by a query-relevance matrix $A\in \{0, 1\}^{m\times n}$ and $s(\vx_i, \vw_q) = \langle \vx_i, \vw_q \rangle$ is the inner product\footnote{To present all top-$k$ subset queries, the query-relevance matrix A  is in the huge space $\{0,1\}^{\binom{m}{k} \times m}$.}. They connected the minimal dimension $d$ with the sign-rank from the query-relevance matrix $A$, specifically, ${\rm rank}_{\pm} (2A-1) - 1 \leq d \leq {\rm rank}_{\pm} (2A-1) $, where the sign-rank of a matrix $M\in \real^{m\times n}$ is defined as ${\rm rank}_{\pm} M = \min\{{\rm rank} B | B\in \real^{m\times n} \text{ such that for all } i, j \text{ we have } {\rm sign}B_{ij} = M_{ij}\}$. However, they failed to explicitly construct the relationship in terms of $m$ due to the hardness of computing the sign rank.
Their simulation, on the other hand, employs a \textit{free embedding optimization} to check whether there exists an embeddable configuration of $m$ elements and $\binom{m}{k}$ subsets in $\real^d$. They empirically established a \textit{polynomial} relationship between $d$ and $m$. \textbf{It is this empirical fitting result that suggested the minimal dimension $d$ required grows with the number of elements $m$ to be retrieved, and indicated the fundamental limit lies in the geometric space.}


% Please add the following required packages to your document preamble:
\begin{table*}[]
\caption{Tight bounds of minimal dimensions in the standard setting (MED) and the centroid setting (MED-C) of important scoring functions. Theorems are stated in Section~\ref{sec:optimal-standard-theory} and Section~\ref{sec:maed_and_exp}. $m$ denotes the size of the universe, $k$ indicates the largest size subset to query. Interestingly, the MED does not depend on $m$. The $O(\log m)$ upper bound of MED-C can be easily observed in simulation.MED is the lower bound of MED-C by definition, so the $\Omega(k)$ lower bound of MED-C is omitted in the table.}\label{tab:major_results}
\centering
% \resizebox{\linewidth}{!}{
\begin{tabular}{lrrr}
\toprule
\multirow{2}{*}{Scoring function} & \multicolumn{2}{l}{Standard setting (MED)} & Centroid setting (MED-C) \\\cmidrule{2-4}
                                  & Lower bound       & Upper bound      & Upper bound                             \\\midrule
Linear (inner product)            & $k-1$             & $2k$             & $O(k^2\log m)$                             \\
Cosine similarity                 & $k-1$             & $2k+1$           & $O(k^2\log m)$                              \\
Euclidean distance ($\ell_2$)     & $k-1$             & $2k$             & $O(k^2\log m)$                            \\\bottomrule
\end{tabular}
% }
\end{table*}


This paper questions the conclusions by \citet{weller2025theoretical} by studying the MED problem, as will be formalized in Section~\ref{sec:problem-definition}, from the perspective of approximability. That means we also work on the fundamental property of geometric spaces as in \citet{weller2025theoretical} but with an extended set of essential scoring functions: in addition to the inner product previously discussed~\citep{weller2025theoretical}, cosine similarity and Euclidean distances are also included due to their wide application in embedding-based systems and representation learning. This paper discusses \textbf{when those spaces are large enough} to contain the subset membership of at most $k$ elements. This perspective, due to its theoretical nature, is not practical enough to help develop scalable ML algorithms. However, \textit{it still allows people to determine whether the effectiveness of embedding-based systems is limited by} \textbf{approximability} (a fundamental property of geometric spaces) \textit{or by} \textbf{learnability} (the way we build specific retrieval systems). If it turns out to be an approximability issue, any methodological effort might make little progress before the dimension of the vector space is large enough. Otherwise, there is still hope in developing more performant embedding-based retrieval systems.


\textbf{Theoretical Contribution.} We establish a clear quantitative relation $d = \Theta(k)$ between the MED $d$ and the maximum cardinality $k$ of the answer set for all three scoring functions. The major results are presented in Table~\ref{tab:major_results}, ``Standard setting (MED)''. Our theoretical results are independent of the universe's size $m$, suggesting that the fundamental property of geometric space \textbf{does not} limit the embedding-based retrieval system. They sharpen the theoretical bounds by~\citet{weller2025theoretical} and contradict the empirical results in~\citet{weller2025theoretical}.

\textbf{Empirical Support.} We also conduct numerical simulations to show that the minimal dimension, where an embeddable configuration of embeddings can be found by optimization, does not grow polynomially with $m$, but at most logarithmically. The numerical simulations also justify the correctness of our tight bounds in MED and credit the limitation of the effectiveness of embedding-based systems to learnability rather than approximability. Some brief details are presented here, and more information can be found in Section~\ref{sec:maed_and_exp}.



In our simulation, we \textbf{further assume} the query vector of $q$ to be given by $\vw_q:= \frac{1}{|S|}\sum_{x\in S} \vx$ as the centroid of vectors in the answer set $S$ and denoted as the \textbf{centroid setting}. Our centroid setting is ``more achievable'' than the ``free embedding optimization'' setting~\citep{weller2025theoretical} because our setting only requires optimising $m$ embeddings for elements, but requires no optimization of the $\binom{m}{k}$ embeddings of the subsets. Simulations in both settings only show the existence of an embeddable configuration of $m$ elements in specific dimensions. By definition, they reveal upper bounds on MED.

Another interesting finding is that the simulation in the centroid setting reveals an empirical logarithmic upper bound of MED. This $O(\log m)$ dependency agrees with our theoretical upper bounds from a probabilistic construction. The minimal embeddable dimension in the centroid setting is denoted MED-C, and listed in Table~\ref{tab:major_results}. Compared to the polynomial dependency revealed by the free embedding optimization~\citep{weller2025theoretical}, it might be counterintuitive because the centroid setting has fewer degrees of freedom ($\binom{m}{k}$ subset query embeddings are not optimized in the centroid setting), providing a perfect example of the bottleneck in learnability rather than approximability. More discussion can be found in Section~\ref{sec:maed_and_exp}.

In summary, our work reframes the debate from geometric limits to learnability, offering both theoretical guarantees and empirical evidence that low-dimensional embeddings suffice for perfect retrieval.
This paper is organized as follows. Section~\ref{sec:problem-definition} formally defines the Minimal Embeddable Dimension (MED), establishes some inequalities of MED, and also introduces the centroid setting. Section~\ref{sec:optimal-standard-theory} proves the tight bounds of the MED under various scoring function classes. Section~\ref{sec:maed_and_exp} derives the bounds for MED in centroid settings and validates them through numerical simulations. Section~\ref{sec:related} discusses two highly related works. Finally, Section~\ref{sec:conclusion} concludes the paper with limitations and future directions.
